\documentclass[12pt,a4paper]{article}
\usepackage[utf8]{inputenc}
\usepackage[T1]{fontenc}
\usepackage{amsmath}
\usepackage{booktabs}
\usepackage{geometry}
\usepackage{hyperref}
\usepackage{cite}
\geometry{a4paper,top=2.5cm,bottom=2.5cm,left=2.5cm,right=2.5cm}
\hypersetup{colorlinks=true,linkcolor=blue,citecolor=blue,urlcolor=blue}

\title{\textbf{Heavy Water Electrolysis with Nickel or Palladium Electrodes\\
as a Low-Energy Proton--Electron Capture Initiator:\\
Connection to LENR and the Fleischmann--Pons Effect}}
\author{Eymen Aydoğan}
\date{May 19, 2026}

\begin{document}
\maketitle

\begin{abstract}
We propose a theoretical connection between the Fleischmann--Pons
electrolysis experiment and the PECCNAL proton--electron capture
mechanism. Electrolysis of heavy water (D$_2$O) using nickel or
palladium electrodes drives deuterium ions (D$^+$) into the metal
lattice, where geometric confinement reduces the proton--electron
separation to femtometre scales. We argue this enables
$p + e^- \rightarrow n$ reactions without external energy input,
as the metal crystal provides the required alignment. Produced
neutrons initiate the PECCNAL chain, generating excess heat
detectable via precision calorimetry. This framework provides
a theoretical basis for the anomalous heat observed in LENR
experiments.
\end{abstract}

\noindent\textbf{Keywords:} LENR, cold fusion, heavy water,
electrolysis, palladium, nickel, Fleischmann-Pons, calorimetry

\hrule\vspace{1em}

\section{Introduction}
In 1989, Fleischmann and Pons~\cite{fleischmann1989} reported
anomalous excess heat during electrolysis of D$_2$O using a
palladium cathode. The result was controversial and poorly
reproducible, yet some laboratories reported consistent excess
heat. The Widom--Larsen theory~\cite{widom2006} proposed that
surface proton--electron capture ($p + e^- \rightarrow n$) could
explain LENR phenomena.

We extend this framework using the PECCNAL mechanism and propose
that nickel — far cheaper than palladium — may serve as an
effective electrode material.

\section{Metal Lattice as Geometric Aligner}

In a metal crystal, deuterium atoms occupy interstitial sites
separated by $\sim$0.1--0.3 nm. The metal lattice:
\begin{enumerate}
\item Confines D atoms at fixed positions
\item Reduces D--D and D--e$^-$ distances to near-nuclear scales
\item Provides conduction electrons as capture targets
\item Eliminates the need for external proton acceleration
\end{enumerate}

The conduction electron density in nickel is:
\begin{equation}
n_e = \frac{\rho N_A Z_\text{valence}}{M}
= \frac{8908 \times 6.022\times10^{23} \times 2}{58.69}
= 1.83 \times 10^{29}\,\text{m}^{-3}
\end{equation}

This is $10^{12}$ times higher than hydrogen gas at STP,
dramatically increasing capture probability.

\section{Electrolysis Setup}

\begin{table}[h!]
\centering
\caption{Proposed electrolysis configuration}
\vspace{0.5em}
\begin{tabular}{ll}
\toprule
\textbf{Parameter} & \textbf{Value} \\
\midrule
Cathode & Nickel or palladium plate/wire \\
Anode & Stainless steel \\
Electrolyte & D$_2$O + dissolved Li$_2$CO$_3$ \\
Voltage & 3--9 V \\
Current & 0.1--2 A \\
Temperature & Room temperature \\
\bottomrule
\end{tabular}
\end{table}

Li$_2$CO$_3$ serves dual purpose: electrolyte (conductivity)
and Li-6 source (neutron capture for energy extraction).

\section{Reaction Sequence}

\begin{align}
\text{D}^+ + e^-_\text{(Ni)} &\rightarrow n + \nu_e \tag{E1}\\
n + {}^2\text{H} &\rightarrow {}^3\text{H} + p + 2.22\,\text{MeV} \tag{E2}\\
n + {}^6\text{Li} &\rightarrow {}^4\text{He} + {}^3\text{H} + 4.78\,\text{MeV} \tag{E3}
\end{align}

\section{Experimental Signature}

Excess heat detection via calorimetry:
\begin{equation}
Q_\text{excess} = Q_\text{measured} - Q_\text{electrolysis}
= mc\Delta T - VIt
\end{equation}

If $Q_\text{excess} > 0$ and reproducible across multiple trials
with D$_2$O but not H$_2$O, this constitutes evidence for
isotope-dependent nuclear reactions consistent with PECCNAL.

\section{Nickel vs Palladium}

\begin{table}[h!]
\centering
\caption{Electrode material comparison}
\vspace{0.5em}
\begin{tabular}{lcc}
\toprule
\textbf{Property} & \textbf{Palladium} & \textbf{Nickel} \\
\midrule
Cost & $\sim$50 \$/g & $\sim$0.02 \$/g \\
D absorption & Excellent & Good \\
Conduction $e^-$ density & High & High \\
LENR reports & Many & Several (Rossi E-Cat) \\
Availability & Limited & Abundant \\
\bottomrule
\end{tabular}
\end{table}

\section{Conclusion}
Heavy water electrolysis with nickel or palladium electrodes
provides a low-energy pathway to initiate PECCNAL reactions
without external proton acceleration. The metal lattice acts as
a geometric aligner, reducing proton--electron separation to
nuclear scales. Dissolved Li$_2$CO$_3$ enables simultaneous
neutron energy extraction. This experiment is reproducible in
a standard laboratory with minimal equipment.

\section*{Acknowledgments}
The author thanks the LENR research community.

\begin{thebibliography}{9}
\bibitem{fleischmann1989} M. Fleischmann and S. Pons, \textit{J. Electroanal. Chem.}, 261, 301, 1989.
\bibitem{widom2006} A. Widom and L. Larsen, \textit{Eur. Phys. J. C}, 46, 107, 2006.
\bibitem{aydogan2026a} E. Aydoğan, ``Coulomb-Barrier-Free Neutron Production,'' arXiv preprint, 2026.
\bibitem{pdg2022} R. L. Workman et al. (PDG), \textit{Prog. Theor. Exp. Phys.}, 083C01, 2022.
\end{thebibliography}
\end{document}
